\section{课程定位：这门课在 Agent 版图中的角色}

\subsection{讲者背景与课程目标}
本讲由 Yann Dubois 主讲，定位是整个 Agentic AI MOOC 中关于 ``训练 LLM 成为 Agent'' 的总览课。课程不是单一算法讲解，而是把训练、评估、系统、成本、部署连成一个统一工程视角。讲者在开场强调，这是一版补录课程，核心目的是把原课堂因技术问题中断的内容完整补齐。

\begin{importantbox}{核心命题：Agent 能力不是单点突破，而是系统协同}
如果只看模型参数量、只看某个 benchmark 排名，常常会误判进展。课程不断重复一个判断：\textbf{决定上限的是数据质量、评估闭环、系统吞吐、推理策略与成本约束的耦合效果}，而不是任何单个模块。
\end{importantbox}

\subsection{从 ``聊天模型'' 到 ``可执行主体''}
课程把 Agent 定义为能在环境中持续交互、调用工具、跨步骤保持目标一致性的模型系统。与传统单轮问答相比，Agent 有三个结构性变化：
\begin{enumerate}
    \item 输出不再是最终文本，而是中间动作序列（tool call / code / API）。
    \item 成功标准不再是 ``看起来对''，而是是否在外部环境中完成目标。
    \item 推理成本由 ``单次回答'' 转为 ``长时 rollout 的累计资源''。
\end{enumerate}

\begin{knowledgebox}{课程的组织逻辑}
从内容编排看，讲者先定义训练流水线，再讨论为什么 ``evaluation is the key''，随后进入成本与系统瓶颈，最后回到 tool use 和 open-ended evaluation 的落地挑战。这种顺序对应真实研发节奏：先把环路搭起来，再优化每个瓶颈。
\end{knowledgebox}

\subsection{本章小结}
本章建立了课程的基本坐标：这不是 ``某个技巧'' 的集合，而是 Agent 训练的端到端工程方法。后续所有技术细节都围绕一个问题展开，即如何把模型能力稳定转化为可验证、可扩展、可部署的 Agent 行为。


